\clearpage
\chapter{Objetivos concretos y metodología de trabajo}

\section{Objetivo general}

El objetivo principal del trabajo es diseñar una metodología para la optimización de carteras basada en la generación de escenarios sintéticos mediante técnicas de inteligencia artificial.

En concreto, la metodología se orienta al análisis de retornos de 30 activos pertenecientes al S\&P 500 y emplea un Autoencoder Variacional como modelo generativo ejecutado. El propósito es estudiar si los escenarios obtenidos desde su espacio latente pueden proporcionar una base alternativa para estimar los parámetros de una cartera. Las arquitecturas VAE-GAN y WGAN pertenecen al marco conceptual y a la prospectiva, no al experimento validado.

Para ello, se plantea comparar los resultados obtenidos mediante la metodología propuesta con carteras construidas a partir del modelo de media-varianza de Markowitz y con una cartera equiponderada. De este modo, se busca evaluar si el uso de modelos generativos puede aportar información adicional en el proceso de asignación de pesos, especialmente en contextos donde los retornos financieros presentan colas pesadas, asimetrías y dependencias no lineales.

\section{Objetivos específicos}

A continuación, se detallan los objetivos específicos que guiarán el desarrollo del trabajo:

\begin{itemize}
    \item Revisar la teoría moderna de carteras y el modelo de media-varianza de Markowitz como base clásica de la optimización de carteras.
    \item Analizar las limitaciones de la varianza y de la matriz de covarianzas como medidas suficientes para representar el riesgo financiero.
    \item Estudiar conceptos relevantes de la teoría de la información, como la entropía, la divergencia de Kullback-Leibler y la información mutua, en relación con el análisis de distribuciones de retornos.
    \item Revisar los fundamentos de los modelos generativos profundos, especialmente VAE, GAN, WGAN y modelos de difusión.
    \item Construir un pipeline de adquisición de datos financieros de activos pertenecientes al S\&P 500.
    \item Preprocesar los datos financieros para su uso posterior, incluyendo el tratamiento de precios ajustados, cálculo de log-retornos, limpieza de datos y normalización.
    \item Diseñar una arquitectura VAE para obtener una representación latente de la distribución de retornos financieros.
    \item Situar GAN, WGAN y modelos de difusión en el estado del arte, identificando su posible incorporación en futuras ampliaciones.
    \item Evaluar la calidad de los escenarios generados mediante métricas estadísticas e informacionales.
    \item Comparar las carteras obtenidas mediante la metodología propuesta con carteras construidas mediante el modelo de Markowitz, una cartera equiponderada y el índice S\&P 500 como benchmark.
\end{itemize}

\section{Metodología del trabajo}

\subsection{Alcance y delimitación del trabajo}

Este trabajo tiene un carácter académico y experimental. No pretende ofrecer una recomendación de inversión ni una propuesta de asignación de capital a valores concretos. Tampoco busca predecir de forma puntual los precios futuros de determinados activos financieros. Su objetivo se centra en estudiar la generación de escenarios sintéticos de retornos y su posible utilización dentro de un proceso de optimización de carteras.

El ámbito de aplicación se limita a un conjunto de activos financieros pertenecientes al S\&P 500. La selección final de activos dependerá de la disponibilidad, calidad y completitud de los datos históricos, así como de la capacidad de procesamiento disponible para entrenar los modelos propuestos.

Asimismo, el trabajo estará condicionado por la frecuencia temporal de los datos utilizados y por las decisiones adoptadas en el preprocesamiento, tales como el cálculo de retornos, el tratamiento de valores ausentes, la normalización y la división entre conjuntos de entrenamiento y evaluación.

Los resultados obtenidos deberán interpretarse como evidencia experimental dentro del marco definido en este trabajo. En ningún caso deben entenderse como garantía de rendimiento futuro ni como recomendación financiera aplicable directamente a decisiones reales de inversión.

\subsection{Justificación de la metodología CRISP-DM}
CRISP-DM se adopta porque proporciona una estructura iterativa que conecta la pregunta de investigación con la evidencia producida en cada etapa. En un problema financiero no basta con entrenar una red neuronal: las conclusiones dependen de cómo se seleccionan los activos, cómo se alinean las fechas, qué transformaciones se aplican y de qué manera se separa la información disponible de la evaluación futura. La metodología obliga a hacer explícitas estas decisiones y facilita volver a fases anteriores cuando el EDA o la evaluación revelan limitaciones.

Esta adopción sigue el modelo de proceso independiente del sector y de la tecnología descrito por \cite{wirth2000crispdm}.

Su carácter independiente de una herramienta concreta también resulta adecuado para este trabajo. La adquisición se realiza con \texttt{yfinance}, la preparación con pandas y scikit-learn, el modelado con PyTorch y la optimización con SciPy; CRISP-DM permite describir el proceso por objetivos y evidencias, no por librerías. Las fases se materializan en scripts, CSV, figuras y un modelo guardado, de modo que la memoria puede vincular cada afirmación empírica con un artefacto verificable.

Aunque CRISP-DM suele representarse de forma secuencial, aquí se entiende como ciclo. El análisis de los escenarios, por ejemplo, puede conducir a revisar el peso de la regularización o la dimensión latente; del mismo modo, una cartera excesivamente concentrada puede motivar nuevas restricciones. Esta posibilidad de retroalimentación es coherente con el carácter experimental del TFM.

\subsection{Adaptación de CRISP-DM al presente trabajo}

\subsubsection{Comprensión del problema}
La primera fase delimita la decisión que se quiere apoyar. El objetivo no es adivinar el precio de cada acción, sino construir una distribución de posibles retornos conjuntos que permita estimar medias y covarianzas alternativas. El punto de comparación es la sensibilidad de Markowitz a dichos parámetros. Por ello se definen tres referencias: equiponderación, Markowitz con posiciones largas y Markowitz con límite del 20~\% por activo.

Esta fase también fija los criterios de éxito y de prudencia. Una estrategia se evalúa mediante retorno acumulado y anualizado, volatilidad, Sharpe y máximo \textit{drawdown}; ninguna métrica aislada basta para afirmar superioridad. Además, se establece que los pesos deben obtenerse sin consultar el conjunto de test y que el resultado no tiene finalidad de asesoramiento financiero. El Capítulo 5 traduce estas decisiones en funciones de optimización y un protocolo común.

\subsubsection{Comprensión de los datos}
La segunda fase examina la cobertura y la estructura estadística antes de modelar. El conjunto final contiene 2.766 sesiones por 30 activos entre enero de 2014 y diciembre de 2024. Se analizan valores ausentes, duplicados, estadísticos descriptivos, asimetría, curtosis, volatilidad móvil y correlaciones. Los cinco CSV y siete figuras generados por \texttt{generate\_eda\_outputs.py} constituyen la evidencia reproducible de esta fase.

El EDA permite comprobar que la matriz final no contiene ausencias y que los retornos presentan valores extremos, curtosis positiva y dependencia transversal. Estos hallazgos justifican estudiar un modelo multivariante flexible, pero no anticipan que vaya a mejorar la cartera. La interpretación detallada se reserva para la Sección 5.1.

\subsubsection{Preparación de los datos}
Los cierres ajustados se ordenan cronológicamente y se filtran con un umbral máximo de ausencias del 5~\%. Los huecos puntuales se propagan hacia delante y después se calculan log-retornos. La transformación elimina la primera fila indefinida y produce una matriz fecha--activo alineada. Se trabaja con retornos, en lugar de precios, para reducir la dependencia del nivel nominal y comparar movimientos relativos.

La división es temporal: las primeras 2.212 observaciones forman entrenamiento y las 554 finales, test. El escalador Min--Max se ajusta exclusivamente con entrenamiento y transforma ambos subconjuntos a $[-1,1]$. Guardar el escalador permite deshacer la transformación de los escenarios. Este orden evita que mínimos, máximos o retornos del periodo de evaluación influyan en el entrenamiento.

\subsubsection{Modelado}
La implementación principal es un VAE con entrada de 30 variables, dos capas ocultas de 64 unidades, espacio latente de dimensión ocho y salida Tanh. Se entrena con Adam durante 150 épocas, lotes de 64 y una pérdida formada por MSE más $0{,}001$ veces la divergencia KL. La semilla 42 favorece la repetibilidad.

Para responder al requisito de comparación se ejecuta además un experimento homogéneo de 120 épocas: autoencoder determinista con latente ocho, VAE base con latente ocho y VAE con latente cuatro. Los tres usan el mismo \textit{split}, capas ocultas, tasa de aprendizaje y semilla. La comparación distingue reconstrucción y calidad estadística de las muestras; no confunde estas métricas con el rendimiento de cartera.

\subsubsection{Evaluación}
La evaluación se organiza en tres niveles. En reconstrucción se calculan MSE y MAE para entrenamiento y test. En generación se comparan errores absolutos de medias y volatilidades, y el error medio de las matrices de correlación. Finalmente, las estrategias se comparan sobre los mismos retornos reales de test mediante las cinco métricas financieras definidas previamente.

La separación de niveles es necesaria porque un autoencoder puede reconstruir bien sin generar muestras útiles, mientras un modelo más regularizado puede perder precisión puntual y producir un espacio latente más muestreable. Del mismo modo, una mayor rentabilidad puede venir acompañada de volatilidad y caída superiores. La discusión del Capítulo 5 conserva estas tensiones.

\subsubsection{Despliegue y reproducibilidad}
En este trabajo el despliegue no consiste en publicar un servicio de inversión, sino en dejar un pipeline repetible. Los scripts se organizan por datos, análisis, modelos, optimización y evaluación. Las entradas y salidas se almacenan en rutas explícitas; los resultados numéricos quedan en CSV y las visualizaciones en PNG. \texttt{requirements.txt} documenta las dependencias y las semillas se fijan cuando existe aleatoriedad.

La trazabilidad permite reconstruir la cadena desde \texttt{sp500\_prices\_adj\_close.csv} hasta \texttt{benchmark\_comparison.csv}. También se conservan el \textit{checkpoint} del VAE, el escalador, las reconstrucciones y los escenarios. El Capítulo 6 describe el orden de ejecución y las limitaciones para reproducir exactamente una descarga futura de un proveedor externo.

\begin{table}[htbp]
\centering
\small
\begin{tabular}{p{0.23\textwidth}p{0.39\textwidth}p{0.25\textwidth}}
\toprule
Fase CRISP-DM & Aplicación en el TFM & Evidencia generada \\
\midrule
Comprensión del problema & Pregunta, benchmarks y métricas & Formulación y restricciones \\
Comprensión de los datos & Cobertura, calidad y EDA & Cinco CSV y siete figuras \\
Preparación de datos & Retornos, corte temporal y escalado & Ficheros train/test y escalador \\
Modelado & AE, VAE y escenarios & Modelos, pérdidas y CSV comparativo \\
Evaluación & Calidad estadística y carteras & Métricas y gráficos \\
Reproducibilidad & Encadenamiento de scripts & Repositorio y artefactos \\
\bottomrule
\end{tabular}
\caption{Adaptación de CRISP-DM al desarrollo experimental del trabajo. Fuente: elaboración propia.}
\label{tab:crisp_dm_tfm}
\end{table}

\subsection{Visión general y secuencia experimental}
El diseño completo encadena adquisición, retornos, EDA, partición, escalado, entrenamiento, generación, estimación sintética, optimización y evaluación. Las decisiones de cartera se toman con entrenamiento histórico o escenarios derivados de él; test solo se utiliza al final. Este diseño ofrece una comparación fuera de muestra coherente, aunque un único corte no sustituye una validación \textit{walk-forward}.

\clearpage
\chapter{Metodología}

\section{Diseño general del experimento}

En este capítulo se presenta la metodología de trabajo seguida para la realización del presente trabajo de investigación. Se detalla el diseño experimental y la lógica de construcción de la arquitectura híbrida VAE-GAN. Asimismo, se describe el proceso de adquisición y preprocesamiento de datos, la formalización matemática de las funciones de pérdida y el protocolo de validación para la selección de carteras.

La metodología propuesta se estructura en varias etapas. En primer lugar, se adquieren datos históricos de activos pertenecientes al S\&P 500 mediante la herramienta \textit{yfinance}. En segundo lugar, se realiza el preprocesamiento de los datos para obtener una representación adecuada de los retornos financieros. Posteriormente, se emplea una arquitectura híbrida VAE-GAN para extraer una representación latente de los datos y generar escenarios sintéticos de mercado. Finalmente, se plantea la optimización de la cartera y la comparación con distintos modelos de referencia.

\section{Adquisición de datos}

Una parte fundamental del experimento propuesto es la adquisición de datos reales que permitan llevar a cabo un estudio relevante.

Empleando la herramienta \textit{yfinance}, se creará un pipeline de datos para adquirir series históricas de los últimos 10 años de los principales activos del S\&P 500.

Se emplearán estos datos por su relevancia. El mercado estadounidense es uno de los más importantes en cuanto a capitalización de mercado. Este índice agrupa a 500 de las mayores empresas de Estados Unidos y constituye una referencia ampliamente utilizada en el análisis financiero. Por su diversificación y representatividad, se considera un escenario adecuado para estudiar metodologías de optimización de carteras.

Asimismo, se procurará establecer la trazabilidad de los datos, con el fin de asegurar que proceden de fuentes confiables y que el proceso de adquisición pueda ser reproducible.

\section{Preprocesamiento de los datos}

Tras la adquisición de los datos, se llevarán a cabo las transformaciones necesarias para poder desarrollar los objetivos del trabajo. Este preprocesamiento permitirá adaptar las series históricas al formato requerido por los modelos generativos y por el proceso posterior de optimización.

En primer lugar, se trabajará con precios ajustados, de forma que se tengan en cuenta eventos corporativos como dividendos o desdoblamientos de acciones. A partir de estas series se calcularán los retornos financieros, que constituirán la entrada principal del modelo.

Para el preprocesado, se aplicará una normalización Min-Max tras un tratamiento de valores atípicos (\textit{outliers}). Este escalado al rango $[0,1]$ resulta relevante para asegurar la estabilidad numérica de las funciones de activación de las redes neuronales y reducir el riesgo de inestabilidad durante el entrenamiento del VAE y la GAN.

Asimismo, se realizará una revisión de valores ausentes o inconsistentes, así como las transformaciones adicionales que sean necesarias para garantizar que los datos puedan ser utilizados de forma adecuada dentro del pipeline experimental.

\section{Arquitectura VAE-GAN propuesta}

Una vez extraídos y preprocesados los datos, se elaborará una estructura en línea con lo expuesto en los apartados previos.

La arquitectura propuesta para el modelo se articula en tres módulos principales:

\begin{itemize}
\item Extracción de características con Autoencoders Variacionales (VAEs) para la reducción de la dimensionalidad de forma no lineal. El VAE proyecta los datos históricos hacia un espacio latente de menor dimensión. Este proceso permite reducir parte del ruido, reteniendo las características persistentes de los datos.

```
\item Generación de escenarios mediante Redes Generativas Adversarias (GANs). La GAN entrena sobre el espacio latente y sintetiza vectores de retornos orientados a preservar dependencias no lineales y momentos estadísticos de orden superior.

\item Optimización de la cartera mediante el uso de medidas informacionales, como la Divergencia de Kullback-Leibler, como criterio complementario al análisis clásico del riesgo. De este modo, se plantea evaluar la robustez de la cartera frente a escenarios generados por el modelo.
```

\end{itemize}

De manera intuitiva, el modelo emplea el VAE para obtener una representación latente de los datos, incorporando información sobre relaciones no lineales entre los activos. A partir de dicho espacio latente, la GAN sintetiza escenarios de mercado que permiten analizar la robustez de distintas asignaciones de cartera.

El mercado financiero presenta una gran cantidad de variables interconectadas de formas que van más allá de la linealidad. Por ello, el VAE se utiliza para aproximar un espacio latente en el que se representen relaciones subyacentes entre los activos. Este planteamiento resulta adecuado frente a métodos estrictamente lineales como el PCA, dado que los mercados financieros no presentan necesariamente una estructura lineal.

El espacio latente $Z$ representa, por tanto, una codificación comprimida de la información relevante contenida en los datos de entrada.

\begin{figure}[H]
\centering
\shorthandoff{>}
\begin{tikzpicture}[
node distance=1.2cm and 1cm,
every node/.style={draw, rectangle, rounded corners, align=center, minimum width=2.2cm, minimum height=0.9cm, font=\footnotesize},
arrow/.style={->, >=Stealth, thick}
]

```
% --- NIVEL 1: PROCESO VAE (Superior) ---
\node (data) {Datos\\ históricos $x$};
\node (enc) [right=of data] {Encoder\\ $q_{\phi}(z|x)$};
\node (latent) [right=of enc] {Espacio\\ Latente $Z$};
\node (dec) [right=of latent] {Decoder\\ $p_{\theta}(x|z)$};

% --- NIVEL 2: MOTOR GENERATIVO (Centro) ---
\node (gen) [below=of latent, yshift=-0.5cm] {Generador\\ $G(z, \epsilon)$};
\node (noise) [left=of gen] {Ruido $\epsilon$};
\node (disc) [below=of gen, yshift=0.3cm] {Discriminador\\ $D(x)$};

% --- NIVEL 3: RESULTADOS (Derecha) ---
\node (scenarios) [right=of gen, xshift=0.5cm] {Escenarios\\ Sintéticos};
\node (kld) [right=of scenarios] {Cálculo\\ KLD};
\node (opt) [below=of kld, yshift=0.3cm] {Optimización\\ Pesos $w$};

% --- CONEXIONES ---
\draw[arrow] (data) -- (enc);
\draw[arrow] (enc) -- (latent);
\draw[arrow] (latent) -- (dec);

% Conexiones Generación
\draw[arrow] (latent) -- (gen);
\draw[arrow] (noise) -- (gen);
\draw[arrow] (gen) -- (disc);
\draw[arrow] (disc.west) -- ++(-0.5,0) |- node[pos=0.25, left] {feedback} (gen.west);

% Salida a Optimización
\draw[arrow] (gen) -- (scenarios);
\draw[arrow] (scenarios) -- (kld);
\draw[arrow] (kld) -- (opt);
```

\end{tikzpicture}
\shorthandon{>}
\caption{Arquitectura híbrida VAE-GAN para la generación y optimización de carteras.}
\end{figure}

El VAE toma datos históricos de retornos reales para obtener variables latentes asociadas a los mismos, modelando relaciones subyacentes entre distintos activos. El Generador recibe este espacio latente y, trabajando junto con el Discriminador, introduce variaciones mediante ruido aleatorio. Con este modelo, se busca crear escenarios sintéticos suficientemente similares a los retornos históricos para que resulten útiles en el proceso de optimización. Esto permite realizar la optimización final sobre un conjunto ampliado de escenarios probables, generando un conjunto de pesos para la cartera.

\section{Generación de escenarios sintéticos}

Una vez estructurado el espacio latente, se procede a la generación de escenarios sintéticos. Para ello, se emplea la arquitectura adversarial de las Redes Generativas Adversarias. En este enfoque, se introduce un marco de entrenamiento adversario en el que dos redes neuronales compiten entre sí: el Generador ($G$) y el Discriminador ($D$).

El Generador se encarga de generar escenarios sintéticos que sean indistinguibles de los datos reales, mientras que el Discriminador se encarga de distinguir entre los datos reales y los sintéticos. El Generador se entrena para maximizar la probabilidad de error del Discriminador, mientras que este último se entrena para maximizar la probabilidad de distinguir entre los datos reales y los sintéticos.

El Generador entrena para buscar escenarios que sean indistinguibles de la realidad, utilizando la estructura latente previamente establecida, así como ruido aleatorio introducido en la entrada.

Por lo tanto, se plantea un juego minimax de suma cero donde el objetivo es alcanzar un equilibrio de Nash. En este escenario, el Generador ($G$) busca maximizar la probabilidad de error del Discriminador ($D$), mientras que este último se optimiza para distinguir con precisión entre la distribución de datos reales $p_{data}$ y la distribución sintética $p_g$.

\section{Función de pérdida y criterios de optimización}

Siguiendo el enfoque propuesto por \cite{kingma2013auto}, el objetivo del modelo VAE es maximizar la verosimilitud marginal de los datos observados. Dado un conjunto de datos ${x^{(i)}}_{i=1}^{N}$, la verosimilitud marginal se descompone como la suma sobre las observaciones individuales:

\begin{equation}
\log p_{\theta}(x^{(1)}, \dots, x^{(N)}) =
\sum_{i=1}^{N} \log p_{\theta}(x^{(i)})
\end{equation}

Sin embargo, el cálculo directo de $\log p_{\theta}(x^{(i)})$ resulta intratable debido a la integración sobre las variables latentes. Para abordar este problema, se introduce una distribución aproximada $q_{\phi}(z|x)$, lo que permite descomponer la log-verosimilitud como:

\begin{equation}
\log p_{\theta}(x^{(i)})
=
D_{KL}
\left(
q_{\phi}(z|x^{(i)}) ,|, p_{\theta}(z|x^{(i)})
\right)
+
\mathcal{L}(\theta, \phi; x^{(i)})
\end{equation}

Dado que la divergencia de Kullback-Leibler es siempre no negativa, el término $\mathcal{L}(\theta, \phi; x^{(i)})$ constituye una cota inferior variacional (Evidence Lower Bound, ELBO):

\begin{equation}
\log p_{\theta}(x^{(i)}) \geq \mathcal{L}(\theta, \phi; x^{(i)})
\end{equation}

Esta cota puede expresarse como:

\begin{equation}
\mathcal{L}(\theta, \phi; x^{(i)})
=
\mathbb{E}*{q*{\phi}(z|x^{(i)})}
\left[
\log p_{\theta}(x^{(i)}, z) - \log q_{\phi}(z|x^{(i)})
\right]
\end{equation}

o, de forma equivalente, separando sus componentes principales:

\begin{equation}
\mathcal{L}(\theta, \phi; x^{(i)})
=
-D_{KL}
\left(
q_{\phi}(z|x^{(i)}) ,|, p(z)
\right)
+
\mathbb{E}*{q*{\phi}(z|x^{(i)})}
\left[
\log p_{\theta}(x^{(i)}|z)
\right]
\end{equation}

Esta formulación revela dos términos fundamentales:

\begin{itemize}
\item Un término de regularización, que fuerza la distribución latente aproximada $q_{\phi}(z|x)$ a mantenerse cercana a la probabilidad a priori $p(z)$.
\item Un término de reconstrucción, que mide la capacidad del modelo generativo para reproducir los datos observados a partir de las variables latentes.
\end{itemize}

El proceso de entrenamiento consiste en maximizar esta cota inferior respecto a los parámetros del modelo generativo $\theta$ y los parámetros variacionales $\phi$.

No obstante, la optimización presenta dificultades prácticas. En particular, el cálculo del gradiente respecto a $\phi$ mediante estimadores Monte Carlo directos presenta alta varianza, lo que dificulta la convergencia y hace inviable su uso en la práctica. Para superar esta limitación, se introduce el truco de reparametrización, que permite expresar la variable latente como una transformación determinista de una variable aleatoria auxiliar:

\begin{equation}
z = \mu_{\phi}(x) + \sigma_{\phi}(x) \cdot \epsilon,
\quad
\epsilon \sim \mathcal{N}(0, I)
\end{equation}

Esta reformulación permite trasladar la fuente de aleatoriedad fuera de la red neuronal, haciendo posible aplicar técnicas estándar de optimización basadas en gradiente (\textit{backpropagation}) con estimadores de baja varianza.

En conjunto, este marco convierte el problema de inferencia variacional en uno de optimización eficiente, permitiendo el entrenamiento conjunto del encoder y el decoder dentro de una arquitectura diferenciable end-to-end.

\section{Optimización y selección de pesos}

Tras la generación de escenarios sintéticos, se procede a la fase de optimización y selección de pesos de la cartera. El objetivo de esta etapa es obtener un vector de pesos $w$ que permita construir una cartera robusta bajo los escenarios generados por el modelo.

La optimización final se realiza sobre escenarios sintéticos generados a partir de la estructura latente aprendida. De este modo, la selección de pesos no depende únicamente de la muestra histórica observada, sino también de escenarios alternativos que preservan ciertas propiedades estadísticas de los retornos reales.

En esta fase se plantea incorporar medidas informacionales, como la Divergencia de Kullback-Leibler, como criterio complementario para evaluar la distancia entre la distribución empírica y la distribución generada. Esta idea permite conectar el proceso de optimización con el marco de teoría de la información desarrollado en capítulos anteriores.

\section{Benchmarks de comparación}

La validación final de la cartera se realizará mediante la comparación con diferentes modelos de elaboración de carteras. En particular, se plantea comparar la metodología propuesta con enfoques clásicos de referencia.

El primer benchmark será el modelo de media-varianza de Markowitz, por constituir la referencia clásica en optimización de carteras. El segundo benchmark será una cartera equiponderada, en la que todos los activos seleccionados reciben el mismo peso. Adicionalmente, se podrá utilizar el índice S\&P 500 como referencia de mercado.

La comparación se realizará atendiendo a métricas financieras y estadísticas que permitan evaluar el comportamiento relativo de cada aproximación. Entre ellas se considerarán el retorno obtenido, la volatilidad, la estabilidad de los pesos y la robustez de la cartera frente a distintos escenarios de mercado.

